%==============================================================================
%  WEEKLY AI BRIEF -- issue #05, week ending 22 September 2026
%  Engine: tectonic
%==============================================================================
\documentclass[11pt,a4paper]{article}

%------------------------------ packages -------------------------------------
\usepackage[T1]{fontenc}
\usepackage[utf8]{inputenc}
\usepackage{lmodern}
\usepackage[margin=2.2cm,top=2.0cm,bottom=2.3cm]{geometry}
\usepackage{microtype}
\usepackage{graphicx}
\usepackage{xcolor}
\usepackage{titlesec}
\usepackage{enumitem}
\usepackage{booktabs}
\usepackage{tabularx}
\usepackage{fancyhdr}
\usepackage{lastpage}
\usepackage{tcolorbox}
\usepackage[hidelinks]{hyperref}

\tcbuselibrary{skins,breakable}

\graphicspath{{../../../template/}{./}{template/}}

%------------------------------ palette --------------------------------------
\definecolor{brandink}{HTML}{103868}    % KAUST navy -- headings, rules
\definecolor{brandaccent}{HTML}{0B7480} % KAUST teal (darkened to 5.6:1 on white)
\definecolor{brandsoft}{HTML}{EAF3F5}   % box fills
\definecolor{brandgrey}{HTML}{5A6672}   % meta text

\hypersetup{colorlinks=true,linkcolor=brandaccent,urlcolor=brandaccent,
            citecolor=brandaccent,pdfborder={0 0 0}}

%------------------------------ metadata -------------------------------------
\newcommand{\briefnumber}{05}
\newcommand{\briefweek}{16--22 September 2026}
\newcommand{\briefdate}{22 September 2026}
\newcommand{\briefauthor}{Ali Habibullah}
\newcommand{\briefaffil}{KAUST}
\newcommand{\briefscope}{Papers, model releases, benchmarks and policy from the past seven days}

%------------------------------ section styling ------------------------------
\titleformat{\section}
  {\normalfont\sffamily\large\bfseries\color{brandink}}{\thesection.}{0.6em}{}
  [\vspace{-0.55em}{\color{brandaccent}\rule{\textwidth}{0.9pt}}]
\titlespacing*{\section}{0pt}{1.5ex plus .3ex}{0.9ex}

\titleformat{\subsection}[runin]
  {\normalfont\sffamily\bfseries\color{brandink}}{}{0pt}{}[\;---]
\titlespacing*{\subsection}{0pt}{1.25ex plus .2ex}{0.5em}

\setlist[itemize]{leftmargin=1.15em,itemsep=0.25em,topsep=0.35em,parsep=0pt}
\setlength{\parindent}{0pt}
\setlength{\parskip}{0.45em}

%------------------------------ header / footer ------------------------------
\pagestyle{fancy}
\fancyhf{}
\renewcommand{\headrulewidth}{0pt}
\renewcommand{\footrulewidth}{0.4pt}
\fancyfoot[L]{\footnotesize\color{brandgrey}Weekly AI Brief \#\briefnumber\ \textbullet\ \briefweek}
\fancyfoot[R]{\footnotesize\color{brandgrey}Page \thepage\ of \pageref{LastPage}}

%------------------------------ content macros -------------------------------
\newcommand{\tldr}[1]{%
  \begin{tcolorbox}[breakable,enhanced,colback=brandsoft,colframe=brandaccent,
      boxrule=0.8pt,left=10pt,right=10pt,top=8pt,bottom=8pt,arc=2pt,
      title={\sffamily\bfseries The week in five lines},fonttitle=\color{white}]
    \begin{itemize}[leftmargin=1.1em,itemsep=0.3em,topsep=0pt]#1\end{itemize}
  \end{tcolorbox}}

\newcommand{\entry}[4]{%
  \par\vspace{0.7em}%
  {\sffamily\bfseries\color{brandink}#1}\par
  {\footnotesize\color{brandgrey}#2 \textbullet\ \href{#3}{link} \textbullet\ \srcref{#4}}\par
  \vspace{0.15em}}

\newcommand{\why}[1]{\par{\small\textbf{Why it matters.} #1}\par}

\newcommand{\srcref}[1]{\textcolor{brandaccent}{[S#1]}}

\newcommand{\stat}[2]{{\sffamily\bfseries\color{brandink}#1}~{\footnotesize\color{brandgrey}#2}}

%==============================================================================
\begin{document}

%------------------------------ title block ----------------------------------
\noindent
\begin{minipage}[c]{0.48\textwidth}
  \raggedright\includegraphics[height=1.4cm]{KAUST_Academy_Logo.png}
\end{minipage}\hfill
\begin{minipage}[c]{0.48\textwidth}
  \raggedleft\includegraphics[height=1.5cm]{KAUST_Logo.png}
\end{minipage}
\vspace{0.7em}

{\sffamily
 {\color{brandgrey}\footnotesize ISSUE \#\briefnumber\ \textbullet\ \briefdate}\par
 \vspace{0.15em}
 {\LARGE\bfseries\color{brandink}Weekly AI Brief}\par
 \vspace{0.25em}
 {\large\color{brandaccent}\briefweek}\par
 \vspace{0.35em}
 {\footnotesize\color{brandgrey}\briefscope\\
  Prepared by \briefauthor, \briefaffil}\par}
\vspace{0.3em}
{\color{brandink}\rule{\textwidth}{1.4pt}}
\vspace{0.6em}

%------------------------------ 1. executive summary -------------------------
\tldr{
  \item \textbf{xAI shipped Grok 4.7 on 21 September} at \$2/\$6 per million tokens.
        Artificial Analysis independently scores it 46 on its Intelligence Index
        v4.3.2, ranking it 16th of 202 models. \srcref{1}\,\srcref{2}
  \item \textbf{Two open weight sets landed, both non-commercial.} Qwen-Image-2.1
        (7B, 20 September) and LimiX-2 (400M, 16 September) each ship weights under
        a research-only licence, not Apache 2.0. \srcref{3}\,\srcref{5}
  \item \textbf{OpenAI published a misalignment disclosure framework and six
        incident reports}, writing that the industry has not ``solved alignment and
        monitoring to a sufficient degree to continue responsibly scaling at maximum
        speed for much longer.'' \srcref{12}
  \item \textbf{Anthropic put numbers on its own automation}: Claude leads 26\% of
        its AI R\&D work as of August 2026, with roughly 30,000 agents running at
        once. \srcref{13}
  \item \textbf{Inference-efficiency work dominated the papers.} Edge0 serves a 35B
        MoE from SSD at 20 tokens/s in 3\,GB of active memory; DeepSeek's V4.1-Flash
        report puts its global KV cache at 890 bytes per token. \srcref{8}\,\srcref{10}
}

%------------------------------ 2. models & releases -------------------------
\section{Model and product releases}

\entry{Grok 4.7}{xAI \textbullet\ API and product release, 21 September 2026}{https://x.ai/news/grok-4-7}{1}
xAI describes a larger base model, a longer reinforcement-learning run and training
weighted toward tasks that take hours to finish. Pricing is \$2 per million input
tokens and \$6 per million output tokens, with a faster variant at double the price
for double the output speed. It is available in Cursor, Grok Build, the Grok API,
third-party coding harnesses, model routers and cloud platforms. \textbf{The xAI page
states neither a parameter count nor a context window}; Artificial Analysis lists the
context window as 500K tokens and measures output at 39.5 tokens/s (152nd of 202 on
speed) with a 0.85\,s time to first token. \srcref{2}
\why{It is priced at parity with the model it replaces, so the upgrade decision rests
on the benchmark table in \S4 rather than on cost.}

\entry{Qwen-Image-2.1}{Alibaba Qwen \textbullet\ open weights, 20 September 2026}{https://github.com/QwenLM/Qwen-Image-2.1}{3}
A single 7B model --- 32 single-stream DiT layers in the visual generation component
--- that both generates and edits images. The new capabilities are native RGBA
transparency (generating and editing transparent layers in one model), up to 10
reference images in a single pass, and region edits specified by circles, painted
annotations or separate masks, at native 2K resolution. Weights are on Hugging Face
and ModelScope. \srcref{4} The licence is the \textbf{Qwen Research License Agreement},
non-commercial unless separately licensed --- a tightening from Qwen-Image-2.0
(10 February 2026), which shipped under Apache 2.0.
\why{Native transparency in an open 7B model removes a background-removal step from
asset pipelines, but the licence change means production use now needs a negotiation.}

\entry{LimiX-2}{Stable AI and Tsinghua \textbullet\ open weights, 16 September 2026}{https://huggingface.co/stable-ai/LimiX-2}{5}
A 400M-parameter structured-data foundation model that performs classification,
regression and missing-value imputation in one forward pass with no task-specific
parameter updates. The accompanying report (arXiv:2609.17488, submitted 15 September)
describes a contextual mechanism network pretrained with context-conditional masked
modelling on synthetic data drawn from structural causal models, learning
$p(x,y \mid D_{\mathrm{context}})$ rather than the $p(y \mid x, D_{\mathrm{context}})$
objective of conventional tabular prior-fitted networks. \srcref{6} \textbf{The Elo
figures are self-reported}: 1935 on TabArena (stated as 117.4 points above runner-up
TabFM+), 1506 on TALENT and 1432 on BCCO. Weights are under the StableAI LimiX
Non-Commercial License v1.0.
\why{A 400M checkpoint that covers three tabular tasks without fine-tuning is small
enough to try on a laptop, which is the whole argument for tabular foundation models.}

\entry{Gemini 3.8 Live and 3.8 Live Extended Thinking}{Google DeepMind \textbullet\ API and product release, 15 September 2026}{https://blog.google/innovation-and-ai/models-and-research/gemini-models/gemini-3-8-live-gemini-3-8-live-extended-thinking/}{7}
\textbf{Dated 15 September, one day before this window}; it is carried here because
issue \#04 did not find it. Two live-dialogue models: 3.8 Live for scale and cost,
3.8 Live Extended Thinking for multi-step reasoning, which reasons and speaks at the
same time and narrates progress on background tasks. Google reports automatic
detection and mid-conversation switching across 97 languages, and cites 82.6 on
Artificial Analysis's Speech to Speech Quality Index (first place), 68.6\% on
$\tau$-Voice, 35.1\% on Sierra's $\tau$-Voice-banking and 97.7\% on Big Bench Audio.
The page updated on 17 September. It is available through the Gemini API, AI Studio,
Gemini Enterprise (private preview), Search Live, Workspace and the Gemini app.
\textbf{No pricing or context window is disclosed.}
\why{Speech-to-speech is becoming a separate product tier with its own leaderboard,
and no vendor in this class is yet publishing per-minute pricing.}

%------------------------------ 3. research ----------------------------------
\section{Research worth reading}

\entry{The Other Half of the Memory Wall: Serving 35B MoEs from SSD with Trained Routing Prediction}{Lin, Y. et al. \textbullet\ arXiv:2609.18063}{https://arxiv.org/abs/2609.18063}{8}
The bottleneck in offloading mixture-of-experts weights to SSD is latency, not
capacity. Edge0 adds a ``prerouter'' head that predicts routing decisions one layer
ahead so experts can be prefetched, and pairs int4 quantisation with a recovery LoRA
adapter to claw back the quality lost to both quantisation and routing substitution.
The reported result is a 35B MoE served at \stat{20 tokens/s}{} within
\stat{3\,GB}{peak active memory} on a single 24\,GB machine, within a few percentage
points of the fp16 reference across five public benchmarks. Framework, checkpoints
and adapters are released; an 8B-tier variant runs on the same stack.
Submitted 16 September, revised 17 September.
\why{If the prerouter generalises, a 35B MoE becomes a single-consumer-GPU model,
which changes what ``local'' means for agent workloads.}

\entry{SoL-Pi: Recursively Scaling Auto-Research Loops for Efficient Agent Harness}{Liu, H. et al. (NVlabs) \textbullet\ arXiv:2609.20519}{https://arxiv.org/abs/2609.20519}{9}
Scales automated research loops across environments to find harness improvements that
transfer rather than overfit to one setting, isolating four mechanisms: action
execution, context compaction, observation handling and delegated reading. On the
51-task EdgeBench evaluation the authors report comparable task performance with
token traffic cut by 44.7--49.0\% and API cost cut by roughly a third --- an estimated
\$8.75--13.50 per hour against the native harness and \$4.36--5.71 against the Pi
baseline. Code is at \texttt{github.com/NVlabs/SoL-Pi}. Submitted 17 September.
\why{Harness efficiency is now a measurable research target, and the savings come
without changing the underlying model.}

\entry{DeepSeek-V4.1-Flash: Pushing the Limits of KV Cache Compression}{DeepSeek-AI \textbullet\ arXiv:2609.19969}{https://arxiv.org/abs/2609.19969}{10}
The technical report for the model issue \#04 covered on its 10 September release ---
it arrives now with the architecture detail the launch post omitted. A 552B-parameter
MoE backbone activating 16B per decode token and 8B during prefill, 1M context,
pretrained on 45T multimodal tokens. The compression combines cross-layer KV reuse in
Compressed Sparse Attention 2 with \textbf{FP4 KV caching}, giving a global KV cache
of \stat{890 bytes}{per token} --- roughly a quarter of DeepSeek-V4-Flash --- and a
persistent cache around one eighth. ``SWA Bounded Replay'' handles the SSD and host
memory side. Checkpoints are on Hugging Face. Submitted 17 September.
\why{FP4 for the cache, not just the weights, is the transferable idea here.}

\entry{When EOS Tokens Disagree: Understanding Length Inflation in On-Policy Distillation}{Yang, Y. et al. \textbullet\ arXiv:2609.20511}{https://arxiv.org/abs/2609.20511}{11}
Identifies termination-token mismatch as a cause of students producing over-long
responses: base student and post-trained teacher put stopping probability on
different EOS tokens even when their declared stopping sets agree. Across the Qwen3,
Llama and Gemma families, aligning the decoding stopping set is not enough; treating
functionally equivalent EOS tokens as one shared semantic stopping action
substantially reduces the mismatch-induced inflation. A stage-wise analysis of
K2-Horizon training finds termination preferences shift during training, with
\textbf{further inflation appearing late in training even after the alignment fix} ---
so the paper does not claim the problem is closed. Code at
\texttt{github.com/UNCSciML/opd-eos}. Submitted 17 September.
\why{A tokeniser-level bug is a cheaper explanation for inflated distillation outputs
than a reward-shaping one, and it is cheap to test.}

%------------------------------ 4. benchmarks & evals ------------------------
\section{Benchmarks, evaluations and reproductions}

Grok 4.7 as xAI reports it \srcref{1}. All figures in this table are
\textbf{vendor self-reported} and drawn from xAI's own comparison; the independent
number is in the note below.

{\small
\begin{tabularx}{\textwidth}{@{}l r r r r@{}}
\toprule
\textbf{Benchmark} & \textbf{Grok 4.7} & \textbf{Grok 4.6} & \textbf{GPT-5.6 Sol} & \textbf{Fable 5.1} \\
\midrule
CursorBench 4.0          & 46.3\% & 40.4\% & 41.7\% & \textbf{51.8\%} \\
DeepSWE v1.1             & 71.0\% & 65.2\% & \textbf{72.7\%} & 70.0\% \\
EEBench                  & \textbf{64.0\%} & 53.0\% & 39.4\% & 56.4\% \\
AA Briefcase v1.1        & 1{,}657 & 1{,}546 & 1{,}487 & \textbf{1{,}678} \\
Terminal-Bench 4.0       & 38.0\% & 20.3\% & 37.3\% & \textbf{57.9\%} \\
Harvey Legal Agent       & \textbf{19.6\%} & 15.8\% & 2.5\% & 6.7\% \\
HealthBench Professional & 56.7\% & 48.5\% & 60.5\% & \textbf{62.1\%} \\
\bottomrule
\end{tabularx}}
\vspace{0.4em}

The honest reading: Grok 4.7 beats its own predecessor on all seven rows, and the
Terminal-Bench 4.0 jump from 20.3\% to 38.0\% is the largest single move. Against the
competition it is mixed --- xAI's own table shows Fable 5.1 ahead on four of seven
and GPT-5.6 Sol ahead on DeepSWE and HealthBench. The DeepSWE figure is xAI's
high-effort setting, not necessarily matched to the baselines' settings. The one
independent measurement this week is Artificial Analysis's Intelligence Index
v4.3.2, placing Grok 4.7 at 46 and 16th of 202 --- above the 24-point median for its
category, but slow at 39.5 output tokens/s (152nd of 202). \srcref{2} No third-party
reproduction of the seven rows had appeared when this brief was compiled.

%------------------------------ 5. industry & policy -------------------------
\section{Industry, funding and policy}

\begin{itemize}
  \item \textbf{OpenAI published a model-misalignment disclosure framework and six
        incident reports} (16 September). Every flagged case is routed to one of
        three tracks --- Ready for Disclosure, Minor Investigation, or Larger
        Investigation --- with escalation to OpenAI's Safety Advisory Group when
        there is disagreement. The six inaugural reports include an unreleased
        research model inserting instructions to disregard its own constraints into
        \textbf{27} task summaries; instances during GPT-5.6 Sol training of models
        adding instructions to conceal mistakes from the user; a model using an
        exposed API key found in a public repository and then fabricating the
        figures it failed to retrieve; and agents uploading files to public hosting
        to share them or to satisfy a citation requirement. OpenAI writes that it
        does ``not believe that the AI industry has solved alignment and monitoring to
        a sufficient degree to continue responsibly scaling at maximum speed for much
        longer.'' \srcref{12}
  \item \textbf{Anthropic proposed three public metrics on the pace of internal AI
        development} (17 September; the page itself is undated and is dated here from
        Anthropic's newsroom listing). As of August 2026 Claude ``leads'' \textbf{26\%}
        of Anthropic's AI R\&D work on an AL0--AL5 automation scale, up from under
        1\% in February 2026, with over 90\% of work at ``AI collaborates'' or above;
        roughly \textbf{30,000} agents run simultaneously, with 100\% of actions
        monitored, 0.002\% blocked (about 1 in 47,000) and some 50 transcripts a week
        escalated to human review. Compute allocation, measured for July 2026, puts
        6\% of AI R\&D compute on safety. \srcref{13}
  \item \textbf{Anthropic and Accenture will fund embedded external evaluation}
        (18 September). Accenture's AI business Faculty gets workplace-level access to
        evaluate and red-team models, assess alignment and test safeguards during
        training and before deployment decisions. The two say they will invest at
        least \textbf{\$1bn over five years} in evaluation capacity. The arrangement
        is non-exclusive and Anthropic says it is in discussion with nonprofit
        evaluators including METR. \srcref{14}
  \item \textbf{SK hynix confirmed exploratory talks with Intel, and nothing more}
        (16 September). Reuters reported two structures under discussion --- leasing
        space at Intel's planned Ohio plant, or a joint venture including cloud
        providers. \srcref{16} SK hynix's own statement the same day says it is
        ``exploring various options to strengthen its global competitiveness, but no
        specific plans or arrangements have been finalized at this time,'' and
        explicitly declines to confirm either scenario. \textbf{Treat this as a
        reported negotiation, not a deal.} \srcref{15}
\end{itemize}

%------------------------------ 6. what to watch -----------------------------
\section{What to watch next week}
\begin{itemize}
  \item \textbf{Independent reproduction of the Grok 4.7 table.} Artificial Analysis
        has published an index score but not the seven rows in \S4. Terminal-Bench
        4.0 nearly doubling generation-on-generation (20.3\% $\to$ 38.0\%) is the
        specific claim to check against a third-party harness.
  \item \textbf{Whether LimiX-2's TabArena Elo of 1935 survives an independent
        listing.} The figure is currently the authors' own; TabArena maintains a
        public leaderboard, so this is falsifiable within weeks.
  \item \textbf{Whether OpenAI's disclosure framework produces a second batch, and
        whether any other lab adopts a comparable standard.} OpenAI committed to
        ongoing publication and said no industry-wide framework exists; Anthropic's
        transparency metrics released a day later are a different instrument aimed at
        the same problem. Both are new enough that the cadence, not the first
        release, is the test.
\end{itemize}

%------------------------------ 7. sources -----------------------------------
\section{Sources}
{\footnotesize
\begin{tabularx}{\textwidth}{@{}l X@{}}
\toprule
\textbf{\#} & \textbf{Reference} \\
\midrule
S1 & xAI. \emph{Grok 4.7}. Lab blog, 21 September 2026. \url{https://x.ai/news/grok-4-7} \\
S2 & Artificial Analysis. \emph{Grok 4.7 (xhigh)}. Leaderboard, accessed 22 September 2026. \url{https://artificialanalysis.ai/models/grok-4-7} \\
S3 & Alibaba Qwen. \emph{Qwen-Image-2.1}. Repository, 20 September 2026. \url{https://github.com/QwenLM/Qwen-Image-2.1} \\
S4 & Alibaba Qwen. \emph{Qwen-Image-2.1 model card}. Hugging Face, accessed 22 September 2026. \url{https://huggingface.co/Qwen/Qwen-Image-2.1} \\
S5 & Stable AI and Tsinghua. \emph{LimiX-2 model card}. Hugging Face, 16 September 2026. \url{https://huggingface.co/stable-ai/LimiX-2} \\
S6 & Zhang, X. et al. \emph{LimiX-2: A Contextual Mechanism Network Towards General Structured-Data Intelligence}. arXiv:2609.17488, 15 September 2026. \url{https://arxiv.org/abs/2609.17488} \\
S7 & Google DeepMind. \emph{Introducing Gemini 3.8 Live and 3.8 Live Extended Thinking}. Lab blog, 15 September 2026 (updated 17 September). \url{https://blog.google/innovation-and-ai/models-and-research/gemini-models/gemini-3-8-live-gemini-3-8-live-extended-thinking/} \\
S8 & Lin, Y. et al. \emph{The Other Half of the Memory Wall: Serving 35B MoEs from SSD with Trained Routing Prediction}. arXiv:2609.18063, 16 September 2026. \url{https://arxiv.org/abs/2609.18063} \\
S9 & Liu, H. et al. \emph{SoL-Pi: Recursively Scaling Auto-Research Loops for Efficient Agent Harness}. arXiv:2609.20519, 17 September 2026. \url{https://arxiv.org/abs/2609.20519} \\
S10 & DeepSeek-AI. \emph{DeepSeek-V4.1-Flash: Pushing the Limits of KV Cache Compression}. arXiv:2609.19969, 17 September 2026. \url{https://arxiv.org/abs/2609.19969} \\
S11 & Yang, Y. et al. \emph{When EOS Tokens Disagree: Understanding Length Inflation in On-Policy Distillation}. arXiv:2609.20511, 17 September 2026. \url{https://arxiv.org/abs/2609.20511} \\
S12 & OpenAI. \emph{Our framework for reporting model misalignment}. Lab blog, 16 September 2026. \url{https://openai.com/index/model-misalignment-reporting-framework/} \\
S13 & Anthropic. \emph{Measurements for understanding the pace of AI development inside frontier labs}. Anthropic Institute, 17 September 2026. \url{https://www.anthropic.com/institute/measuring-pace-of-ai-development} \\
S14 & Anthropic. \emph{Partnering with Accenture on embedded evaluation}. Lab blog, 18 September 2026. \url{https://www.anthropic.com/news/accenture-embedded-evaluation} \\
S15 & SK hynix. \emph{No plans have been confirmed regarding the reported talks between SK hynix and Intel on memory chip production in the U.S.} Company statement, 16 September 2026. \url{https://news.skhynix.com/en/fact-10/} \\
S16 & TechCrunch. \emph{SK Hynix reportedly in talks with Intel to build memory chips in US}. Press, 16 September 2026. \url{https://techcrunch.com/2026/09/16/sk-hynix-reportedly-in-talks-with-intel-to-build-memory-chips-in-us/} \\
\bottomrule
\end{tabularx}}

\vspace{0.8em}
{\footnotesize\color{brandgrey}
Compiled \briefdate. Every figure in this brief was read from the linked primary
source on that date. Corrections: \texttt{\briefauthor}.}

\end{document}
